% ============================================================
% YM-Foundation-A
% A gauge-independent spectral gap from the two-torsion geometry
% of a torus quotient
% ============================================================
\documentclass[11pt,a4paper]{article}
\usepackage{amsmath,amssymb,amsthm}
\usepackage{fontspec}
\setmainfont{Latin Modern Roman}
\newfontfamily\leanmono{DejaVu Sans Mono}[Scale=0.82]
\usepackage{booktabs}
\usepackage{fvextra}
\usepackage{enumitem}
\usepackage[a4paper,margin=2.7cm]{geometry}

% --- Notation ---
\newcommand{\T}{\mathbb{T}}
\newcommand{\R}{\mathbb{R}}
\newcommand{\Z}{\mathbb{Z}}
\newcommand{\C}{\mathbb{C}}
\newcommand{\N}{\mathbb{N}}
\newcommand{\PP}{\mathbb{P}}
\newcommand{\dodd}{d_{\mathrm{odd}}}
\newcommand{\dist}{\mathrm{dist}}
\newcommand{\Hilb}{\mathrm{Hilb}}
\newcommand{\class}{\mathcal{C}}
\newcommand{\thy}{\mathcal{T}}
\newcommand{\Mod}{M}

\theoremstyle{plain}
\newtheorem{theorem}{Theorem}[section]
\newtheorem{lemma}[theorem]{Lemma}
\newtheorem{proposition}[theorem]{Proposition}
\newtheorem{corollary}[theorem]{Corollary}
\theoremstyle{definition}
\newtheorem{definition}[theorem]{Definition}
\theoremstyle{remark}
\newtheorem{remark}[theorem]{Remark}

\title{A gauge-independent spectral gap from the two-torsion geometry\\ of a torus quotient}
\author{Dhiren J. Master}
\date{}

\begin{document}
\maketitle

\begin{abstract}
Let $(M,g)$ be a compact connected Riemannian manifold containing an
embedded copy of the pillowcase, the quotient of the flat two-torus by
the involution $(x,y) \mapsto (-x,-y)$, and let $\Delta$ denote the
minimum geodesic distance from a distinguished corner of the pillowcase
fibre to the remaining three. We prove that the corners are precisely
the four two-torsion points of the torus, pairwise distinct, and that
$\Delta > 0$. We further prove that the framed moduli spaces of
torsion-free sheaves on the projective plane admit a
Bia{\l}ynicki-Birula cellular decomposition with vanishing odd Betti
numbers, so that their Euler characteristics are absolute Betti counts,
and we identify the central Fourier coefficients of the level-two
Appell--Lerch sum with the odd divisor function $\dodd(n)$, with the
coefficient match tabulated to order fifteen. For a class of
Kaluza--Klein-reduced theories defined by three explicit spectral
hypotheses, $\Delta$ is the mass of the lowest non-vacuum one-particle
state, and neither this identification nor the positivity of $\Delta$
refers to the gauge group; the independence is machine-checked.
Existence of the gap therefore transfers across gauge groups, while
its value depends on the corner-selection data. An existence-only
result of this kind may be useful in making progress towards the
mass-gap problem for four-dimensional gauge theories.
\end{abstract}

%%=====================================================================
\section{Introduction}\label{sec:intro}
%%=====================================================================

The purpose of this paper is to establish, from elementary geometric
and arithmetic inputs, the existence of a strictly positive spectral
gap for a class of Kaluza--Klein-reduced theories, and to show that
the existence claim is independent of the gauge group of the reduced
theory. The results divide into an unconditional part and a
conditional part, and the division is maintained explicitly
throughout.

The unconditional part comprises four statements. First, the
fixed-point set of the involution $\sigma(x,y) = (-x,-y)$ on the flat
two-torus $\T^2 = \R^2/\Z^2$ consists of exactly the four two-torsion
points, and these are pairwise distinct; the quotient $\T^2/\sigma$ is
the pillowcase, and the four points are its corners
(Lemma~\ref{lem:corners}). Second, if the pillowcase is embedded as a
fibre of a compact connected Riemannian manifold $(M,g)$ with a
distinguished corner $c$, the minimum geodesic distance $\Delta$ from
$c$ to the remaining three corners is strictly positive
(Lemma~\ref{lem:positivity}). Third, the framed moduli spaces
$\Mod(r,n)$ of torsion-free sheaves on the projective plane admit a
Bia{\l}ynicki-Birula decomposition into affine cells, so that all odd
Betti numbers vanish and the Euler characteristic equals the absolute
Betti count (Theorem~\ref{thm:betti}). Fourth, the central Fourier
coefficients of the level-two Appell--Lerch sum equal the odd divisor
function $\dodd(n)$ (Proposition~\ref{prop:appell}); through the
generating identity of G\"ottsche and the resolution of instanton
moduli by Hilbert schemes of points, coefficients of this kind
organise the counting invariants of the moduli spaces of the third
statement.

The conditional part concerns a class $\class$ of
Kaluza--Klein-reduced theories, defined in Section~\ref{sec:class} by
three explicit spectral hypotheses: the level degeneracies of a theory
in the class equal $\dodd(n)$; the mass at level $n$ equals $n\Delta$;
and $\dodd(n)$ counts one-particle states at level $n$. The hypotheses
are stated as part of the definition of the class and are not derived
here; the unconditional results ground the counting object to which
they refer. For every $\thy \in \class$ we prove that $\Delta$ is the
mass of the lowest non-vacuum one-particle state
(Proposition~\ref{prop:identification}), and that neither this
identification nor the positivity of $\Delta$ makes reference to the
gauge group of $\thy$ (Proposition~\ref{prop:transfer}). The latter
statement is verified mechanically: in the accompanying formal
development the gauge group enters as a parameter that provably does
not occur in the proof term. A corollary delimits the scope: the
existence of the gap is uniform over gauge groups and over admissible
corner selections, whereas the numerical value of $\Delta$ is a
function of the selection data and is not a gauge-group invariant
(Corollary~\ref{cor:scope}).

The method is deliberately elementary. The positivity of $\Delta$
consumes only the two-torsion structure of $\sigma$ and the
Hopf--Rinow theorem. The homological statement rests on the theorem of
Bia{\l}ynicki-Birula together with the torus-action analysis of the
framed moduli spaces due to Nakajima and Yoshioka. The arithmetic
statement is a computation with geometric series occupying half a
page. The conditional results are exact inferences from the stated
hypotheses, with the single internal input $\Delta \ge 0$ discharged
by Lemma~\ref{lem:positivity}. All results in the certified chain,
together with the homological parity implication and the
Appell--Lerch coefficient identity, are accompanied by machine-checked
certificates, reproduced in Appendix~\ref{app:lean}.

The paper is organised as follows. Section~\ref{sec:setting} states
the geometric setting and proves the two lemmas on the corner gap.
Section~\ref{sec:arithmetic} defines the odd divisor function,
establishes its elementary properties, and proves the Appell--Lerch
coefficient identity with a table to order fifteen.
Section~\ref{sec:moduli} proves the vanishing of the odd Betti numbers
of the framed moduli spaces and assembles the generating-function link
between the moduli invariants and the coefficients of
Section~\ref{sec:arithmetic}. Section~\ref{sec:class} defines the
theory class and records a remark on a second, modularity-conditional
route to the absence of a paired sector. Section~\ref{sec:gap} proves
the identification of $\Delta$ with the lowest mass and the
gauge-group independence of the argument. Section~\ref{sec:scope}
delimits the scope of the results and closes with an outlook.
Appendix~\ref{app:lean} reproduces the formal certificates.

%%=====================================================================
\section{The geometric setting and the corner gap}\label{sec:setting}
%%=====================================================================

This section fixes the geometric data and proves the two unconditional
lemmas on which the positivity of the gap rests. Both proofs are
elementary; we work the smallest cases by hand.

\begin{definition}[Setting]\label{def:setting}
Let $\T^2 = \R^2/\Z^2$ be the flat two-torus and let
$\sigma \colon \T^2 \to \T^2$ be the involution
$\sigma(x,y) = (-x,-y)$. The quotient $\T^2/\sigma$ is the
\emph{pillowcase}. Let $(M,g)$ be a compact connected Riemannian
manifold with continuous metric, and let
$\iota \colon \T^2/\sigma \to M$ be a continuous injection of the
pillowcase onto a fibre of $M$. Let $c \in M$ be the image under
$\iota$ of one distinguished corner of the pillowcase (the
\emph{selection}), and define
\begin{equation}\label{eq:gapdef}
  \Delta \;=\; \Delta(g, \iota, c) \;=\;
  \min_{c'} \, \dist_M(c, c'),
\end{equation}
the minimum being taken over the images $c'$ of the three remaining
corners.
\end{definition}

The hypotheses of Definition~\ref{def:setting} are the entire
geometric input of the paper. Which corner is distinguished is data,
not a consequence of the topology; the dependence of the value of
$\Delta$ on this datum is the subject of
Corollary~\ref{cor:scope}.

\begin{lemma}[Corner distinctness]\label{lem:corners}
The fixed-point set of $\sigma$ consists of exactly the four
two-torsion points
\[
\mathrm{Fix}(\sigma)
\;=\;
\bigl\{\, (0,0),\ (\tfrac12, 0),\ (0, \tfrac12),\ (\tfrac12, \tfrac12) \,\bigr\},
\]
and these four points are pairwise distinct in $\T^2$. Consequently
the pillowcase has exactly four corners, and their images under
$\iota$ are four distinct points of $M$.
\end{lemma}

\begin{proof}
We work with representatives $(x,y) \in [0,1)^2$, the fundamental
domain of $\R^2/\Z^2$.

\emph{Fixed points are two-torsion.} A point $[(x,y)]$ is fixed by
$\sigma$ if and only if $(-x,-y) \equiv (x,y) \pmod{\Z^2}$, that is,
if and only if $2x \in \Z$ and $2y \in \Z$. The two conditions are
independent, so it suffices to solve $2x \in \Z$ for $x \in [0,1)$.

\emph{The smallest case, worked by hand.} If $2x = k \in \Z$ with
$x \in [0,1)$, then $k = 2x \in [0,2)$, so $k \in \{0, 1\}$ and
correspondingly $x \in \{0, \tfrac12\}$. Both solve the congruence:
$-0 \equiv 0$ and $-\tfrac12 \equiv \tfrac12 \pmod{\Z}$, their
difference being $-1 \in \Z$. Hence the fixed-point set is exactly
$\{0, \tfrac12\} \times \{0, \tfrac12\}$, the four points displayed.

\emph{Distinctness.} Two representatives in the fundamental domain
$[0,1)^2$ define the same point of $\T^2$ if and only if they are
equal as pairs; the four listed representatives are pairwise unequal,
their coordinatewise differences being $\pm\tfrac12 \notin \Z$ in at
least one coordinate. Hence the four fixed points are pairwise
distinct in $\T^2$. The corners of the pillowcase are by definition
the images of the fixed points in $\T^2/\sigma$, on which the quotient
map is injective; and $\iota$ is injective by hypothesis, so the four
corner images are distinct in $M$.
\end{proof}

We note that Lemma~\ref{lem:corners} is a statement of torus topology:
no structure on the fibre beyond the involution is used, and no
property of any group action other than $\sigma$ enters.

\begin{lemma}[Positivity of the corner gap]\label{lem:positivity}
In the setting of Definition~\ref{def:setting}, each of the three
distances $\dist_M(c, c')$ is strictly positive, and hence
$\Delta > 0$.
\end{lemma}

\begin{proof}
By Lemma~\ref{lem:corners} the four corner images are pairwise
distinct points of $M$. By the Hopf--Rinow theorem
\cite[Ch.~7]{doCarmo}, a compact Riemannian manifold with continuous
metric is a complete metric space in the geodesic distance, and the
distance between two distinct points is strictly positive. Each of
the three distances in \eqref{eq:gapdef} is therefore positive, and
the minimum of three positive numbers is positive. Neither step
invokes any group beyond the involution $\sigma$.
\end{proof}

The formal certificate of Lemma~\ref{lem:positivity}
(Appendix~\ref{app:lean}, Certificate~2) is proved at metric-space
level, which is strictly sufficient for the claim as used; the
Riemannian-geometric source of the metric-space structure is the
Hopf--Rinow theorem as cited.

%%=====================================================================
\section{The odd divisor function and the Appell--Lerch sum}\label{sec:arithmetic}
%%=====================================================================

This section supplies the arithmetic ingredient. We define the odd
divisor function, record the elementary properties consumed later, and
prove that it arises as the sequence of central Fourier coefficients
of the level-two Appell--Lerch sum.

\begin{definition}\label{def:dodd}
For $n \ge 1$, let $\dodd(n)$ denote the number of odd divisors of
$n$. Writing $n = 2^a n_{\mathrm{odd}}$ with $n_{\mathrm{odd}}$ odd,
one has $\dodd(n) = d(n_{\mathrm{odd}})$, the ordinary divisor count
of the odd part.
\end{definition}

\begin{proposition}[Elementary properties]\label{prop:doddprops}
The function $\dodd$ satisfies:
\begin{enumerate}[label=\textup{(\roman*)}]
\item $\dodd(1) = 1$;
\item $\dodd(n) \ge 1$ for all $n \ge 1$;
\item $\dodd(2^a m) = \dodd(m)$ for all $a \ge 0$ and $m \ge 1$;
\item $\dodd(mn) = \dodd(m)\,\dodd(n)$ whenever $\gcd(m,n) = 1$;
\item $\dodd(n) = 1$ if and only if $n$ is a power of two, and
  $\dodd(n) = 2$ if and only if $n_{\mathrm{odd}}$ is prime.
\end{enumerate}
\end{proposition}

\begin{proof}
The sole divisor of $1$ is $1$, which is odd, giving (i); the divisor
$1$ of any $n$ is odd, giving (ii); and (iii) is immediate from
$\dodd(n) = d(n_{\mathrm{odd}})$, since $2^a m$ and $m$ have the same
odd part for odd $m$, and in general the odd part is unchanged by
powers of two. Property (iv) follows from the multiplicativity of $d$
together with the multiplicativity of the odd-part map on coprime
arguments. For (v), $d(n_{\mathrm{odd}}) = 1$ if and only if
$n_{\mathrm{odd}} = 1$, and $d(n_{\mathrm{odd}}) = 2$ if and only if
$n_{\mathrm{odd}}$ is prime.
\end{proof}

\begin{definition}\label{def:appell}
The level-two Appell--Lerch sum is
\begin{equation}\label{eq:appell}
  \mu(\tau,z) \;=\; \sum_{n \ge 1} \frac{q^n}{(1 - yq^n)(1 - y^{-1}q^n)},
  \qquad q = e^{2\pi i\tau},\quad y = e^{2\pi i z},
\end{equation}
regarded as a formal power series in $q$ with coefficients in
Laurent polynomials in $y$.
\end{definition}

Appell--Lerch sums are the canonical building blocks in Zwegers'
theory of mock theta functions \cite{Zwegers}; the specialisation
\eqref{eq:appell} is the level-two member of the family,
corresponding to the order-two quotient in
Definition~\ref{def:setting}. The property we require of $\mu$ is
purely combinatorial.

\begin{proposition}[Central coefficients]\label{prop:appell}
The coefficient of $y^0 q^m$ in $\mu(\tau,z)$ equals $\dodd(m)$ for
every $m \ge 1$.
\end{proposition}

\begin{proof}
Expanding each factor of \eqref{eq:appell} in geometric series,
\[
  \frac{q^n}{(1 - yq^n)(1 - y^{-1}q^n)}
  \;=\; \sum_{a,b \ge 0} y^{\,a-b}\, q^{\,n(1+a+b)} .
\]
The $y^0$ terms are those with $a = b$, so the central part of $\mu$
is
\[
  \sum_{n \ge 1} \sum_{a \ge 0} q^{\,n(2a+1)} .
\]
The coefficient of $q^m$ therefore counts the pairs $(n,a)$ with
$n \ge 1$, $a \ge 0$ and $n(2a+1) = m$. The map $(n,a) \mapsto 2a+1$
is a bijection from this set of pairs onto the set of odd divisors of
$m$, with inverse $d \mapsto (m/d,\, (d-1)/2)$. Hence the coefficient
is $\dodd(m)$.
\end{proof}

The bijection in the proof of Proposition~\ref{prop:appell} is
machine-checked in full generality, and the resulting coefficient
table certified to order fifteen, in Appendix~\ref{app:lean}
(Certificate~7). The certified table is
\begin{center}
\begin{tabular}{c|ccccccccccccccc}
$m$ & 1 & 2 & 3 & 4 & 5 & 6 & 7 & 8 & 9 & 10 & 11 & 12 & 13 & 14 & 15 \\
\midrule
$\dodd(m)$ & 1 & 1 & 2 & 1 & 2 & 2 & 2 & 1 & 3 & 2 & 2 & 2 & 2 & 2 & 4
\end{tabular}
\end{center}
in agreement with Definition~\ref{def:dodd} evaluated directly.

%%=====================================================================
\section{Framed moduli and the vanishing of odd cohomology}\label{sec:moduli}
%%=====================================================================

This section proves the homological statement and assembles the
generating-function context that connects the moduli invariants to
the coefficients of Section~\ref{sec:arithmetic}. Throughout,
$\Mod(r,n)$ denotes the framed moduli space of torsion-free sheaves
$E$ on the projective plane $\PP^2$ with rank $r$ and $c_2(E) = n$,
the framing being a trivialisation of the restriction of $E$ to the
line at infinity $\ell_\infty$; these are the completions of the
moduli of framed instantons on $\C^2 \cong \R^4$, and they are smooth
quasi-projective varieties of dimension $2rn$
\cite{NakajimaYoshioka}.

\begin{theorem}[Vanishing of odd cohomology]\label{thm:betti}
The framed moduli space $\Mod(r,n)$ admits a decomposition into
affine cells of even real dimension. Consequently
$H_{2j+1}(\Mod(r,n), \Z) = 0$ for all $j$, and the Euler
characteristic equals the absolute Betti count:
\begin{equation}\label{eq:chiabs}
  \chi\bigl(\Mod(r,n)\bigr) \;=\; \sum_j (-1)^j\, b_j \;=\; \sum_j b_{2j}
  \;=\; \sum_j b_j .
\end{equation}
There are, in particular, no cancellations: the counting invariant
$\chi(\Mod(r,n))$ is an absolute count of cells.
\end{theorem}

\begin{proof}
The maximal torus of $GL_r(\C)$ acting on the framing, together with
the two-dimensional torus acting on $\PP^2$ preserving
$\ell_\infty$, acts on $\Mod(r,n)$ with finitely many fixed points,
parametrised by $r$-tuples of Young diagrams with $n$ boxes in total
\cite{NakajimaYoshioka}. For a generic one-parameter subgroup of this
torus, the Bia{\l}ynicki-Birula theorem \cite{BB} decomposes
$\Mod(r,n)$ into locally closed attracting cells, one for each fixed
point, each cell being an affine space; the existence of the limits
required by the decomposition in the quasi-projective setting is
established for these moduli spaces in \cite{NakajimaYoshioka} (for
$r = 1$ the statement is the cell decomposition of the Hilbert
schemes of points \cite{NakajimaLectures}; the resulting Betti
numbers of $\Mod(r,n)$ are computed explicitly in
\cite{NYlectures}). Affine cells have even
real dimension, so the cellular chain complex is concentrated in even
degrees, and all odd Betti numbers vanish. The identity
\eqref{eq:chiabs} is then term-by-term: for even $j$ the sign is
$+1$, and for odd $j$ the Betti number is zero. The parity
implication --- vanishing odd part gives signed sum equal to plain
sum, for any Betti sequence --- is machine-checked in full generality
in Appendix~\ref{app:lean} (Certificate~6).
\end{proof}

Theorem~\ref{thm:betti} is the statement that the counting invariants
of the framed moduli spaces are counts without cancellation. Two classical
generating identities organise these counts. First, for $r = 1$ the
moduli space $\Mod(1,n)$ is isomorphic to the Hilbert scheme of $n$
points on $\C^2$ \cite{NakajimaLectures}, and for a smooth projective
surface $S$ G\"ottsche's formula \cite{Gottsche} gives
\begin{equation}\label{eq:gottsche}
  \sum_{n \ge 0} \chi\bigl(\Hilb^n(S)\bigr)\, q^n
  \;=\; \prod_{m \ge 1} (1 - q^m)^{-\chi(S)} .
\end{equation}
The smallest case, worked by hand: $\Hilb^1(S) = S$, and the
coefficient of $q^1$ on the right of \eqref{eq:gottsche} is
$\chi(S)$, as it must be. Second, for higher rank on the projective
plane, the partition functions assembled from the invariants of
$\Mod(r,n)$ in the topologically twisted theory of Vafa and Witten
\cite{VafaWitten} are built from generalised Appell functions: this
is proved by Bringmann, Manschot and Rolen \cite{BMR}, whose
identities express higher-rank counting series on the projective
plane through generalised Appell functions, the family of which the
level-two sum \eqref{eq:appell} is the depth-one member. The central coefficients
of that member are, by Proposition~\ref{prop:appell}, exactly
$\dodd(n)$. The counting object to which the class hypotheses of
Section~\ref{sec:class} refer is therefore assembled entirely from
ingredients each of which stands on the results of this section and
of Section~\ref{sec:arithmetic}: cell counts without cancellation
(Theorem~\ref{thm:betti}), organised by Appell-type series whose
central coefficients are $\dodd(n)$
(Proposition~\ref{prop:appell}).

%%=====================================================================
\section{The theory class}\label{sec:class}
%%=====================================================================

We now define the class of theories for which the gap of
Section~\ref{sec:setting} will be identified with a mass. The
definition is by explicit spectral hypotheses; the class carries the
entire conditional weight of the paper, and nothing in this section
is proved about membership of the class.

\begin{definition}[The class $\class$]\label{def:class}
A \emph{Kaluza--Klein-reduced theory over the setting of
Definition~\ref{def:setting}} is a four-dimensional theory $\thy$
obtained by reduction on a compact internal space containing the
pillowcase fibre, presented through its one-particle spectral data:
a degeneracy function $m \colon \N \to \N$, a mass assignment
$\mathrm{mass} \colon \N \to \R$, and the set
$\mathcal{S}(\thy) \subset \R$ of masses of non-vacuum one-particle
states. The theory $\thy$ belongs to the class $\class$ if:
\begin{enumerate}[label=\textup{(A\arabic*)}]
\item the level degeneracies satisfy $m(n) = \dodd(n)$ for all
  $n \ge 1$;
\item the mass at level $n$ equals $n\Delta$, where $\Delta$ is the
  corner gap \eqref{eq:gapdef} of the internal setting;
\item $\dodd(n)$ counts the one-particle states of $\thy$ at mass
  level $n$; equivalently, $\mathcal{S}(\thy)$ consists exactly of
  the values $\mathrm{mass}(n)$ over the occupied levels $n \ge 1$,
  those with $m(n) \ne 0$.
\end{enumerate}
\end{definition}

Three comments delimit what Definition~\ref{def:class} does and does
not assert. First, the hypotheses (A1)--(A3) are declared, not
derived: whether a given reduction satisfies them is a question about
that reduction, outside the scope of this paper. Second, hypothesis
(A3) is part of the definition of the class and is not a consequence
of Theorem~\ref{thm:betti}. What Theorem~\ref{thm:betti} supplies is
the consistency of the reading: the counting coefficients to which
(A3) refers are absolute counts, free of cancellation, so that a
hypothesis identifying them with numbers of states is a hypothesis
about unsigned quantities. The identification itself remains an
assumption. Third, the deflationary reading of the main results is
therefore the correct one: Section~\ref{sec:gap} proves exact
inferences from (A1)--(A3), and the strength of the conclusions for
any particular theory is precisely the strength of the case that the
theory satisfies the hypotheses.

\begin{remark}[A second route to the absence of a paired sector]
\label{rem:secondroute}
Hypothesis (A3) excludes, by fiat, a paired sector: states cancelling
in a signed count. Theorem~\ref{thm:betti} removes such a sector at
the homological level. It is worth recording that a modular argument
reaches the same exclusion from a different hypothesis. Suppose the
completed counting series of a theory --- the series whose $n$-th
coefficient is the number of states at level $n$, completed so as to
transform as a modular form of weight zero without poles --- differs
from the index series by the generating function of a putative paired
sector, $\sum_n 2k_n q^n$. A holomorphic modular form of weight zero
without poles is constant --- for the full modular group this is
classical \cite[Ch.~VII]{Serre}, and in general it is the statement
that a holomorphic function on a compact Riemann surface is
constant --- and since the vacuum
is unpaired, $k_0 = 0$, the constant is zero, so $k_n = 0$
identically. Weight zero is the natural home for the hypothesis: the
difference of two counting series for the same graded object, with
the same normalisation, carries no automorphy weight of its own, and
constancy is then the full strength of the Liouville-type statement.
The remark is in this sense an arithmetic mirror of
Theorem~\ref{thm:betti}: the homological argument excludes a paired
sector at the level of cells, the modular argument excludes it at the
level of completed series, and neither mechanism implies the other.
The hypothesis of modular completion is not part of the definition of
$\class$ and is not used elsewhere in this paper; the remark records
only that two independent mechanisms both close the paired-sector
loophole under their respective hypotheses.
\end{remark}

%%=====================================================================
\section{The spectral gap and its gauge-group independence}\label{sec:gap}
%%=====================================================================

We now prove the two conditional results. Throughout, $\thy \in
\class$ is a theory in the class of Definition~\ref{def:class}, over
a fixed setting $(M, g, \iota, c)$, and $\Delta$ is the corner gap
\eqref{eq:gapdef}.

\begin{proposition}[Identification of $\Delta$ with the lowest
mass]\label{prop:identification}
Let $\thy \in \class$. Then $\Delta$ is the mass of the lowest
non-vacuum one-particle state of $\thy$; that is, $\Delta$ is the
least element of $\mathcal{S}(\thy)$.
\end{proposition}

\begin{proof}
By (A3) with (A1), the non-vacuum one-particle masses of $\thy$ are
exactly the values $\mathrm{mass}(n)$ over the levels $n \ge 1$ with
$\dodd(n) \ne 0$. The smallest case, worked by hand: $\dodd(1) = 1$,
the sole divisor of $1$ being odd, so level $n = 1$ is occupied. By
(A2), $\mathrm{mass}(n) = n\Delta$; and by
Lemma~\ref{lem:positivity}, $\Delta \ge 0$, so $n\Delta \ge \Delta$
for all $n \ge 1$, with equality at $n = 1$. Hence $\Delta$ belongs
to $\mathcal{S}(\thy)$ and bounds it below.

We record explicitly that the proof consumes the three declared
hypotheses together with one internal input, $\Delta \ge 0$, which is
discharged by Lemma~\ref{lem:positivity} and is not an additional
assumption. The formal certificate (Appendix~\ref{app:lean},
Certificate~3) carries this input as an explicit argument, discharged
from the geometry in Certificate~4, so that the dependence of the
identification on the positivity lemma is visible rather than
silent.
\end{proof}

\begin{proposition}[Gauge-group independence of the corner
gap]\label{prop:transfer}
Let $\thy \in \class$ have gauge group $G$. Then:
\begin{enumerate}[label=\textup{(\roman*)}]
\item $\Delta > 0$;
\item $\Delta$ is the mass of the lowest non-vacuum one-particle
  state of $\thy$;
\item neither \textup{(i)} nor \textup{(ii)} refers to $G$.
\end{enumerate}
Consequently the existence of a positive mass gap holds for every
theory in $\class$, uniformly in the gauge group.
\end{proposition}

\begin{proof}
Claim (i) is Lemma~\ref{lem:positivity}, whose sole geometric input
is Lemma~\ref{lem:corners}. Claim (ii) is
Proposition~\ref{prop:identification} under the hypotheses
(A1)--(A3) of Definition~\ref{def:class}. Claim (iii) is an
inspection of dependencies. The complete input set of (i) and (ii)
consists of five items: the two-torsion structure of $\sigma$ on
$\T^2$ (Lemma~\ref{lem:corners}); the Hopf--Rinow theorem
(Lemma~\ref{lem:positivity}); and the three hypotheses (A1), (A2),
(A3). None of these five inputs mentions $G$: the first two are
statements about the internal geometry, and the three hypotheses
constrain spectral data in which the gauge group does not appear.
The gauge group of $\thy$ is whatever the reduction produces; it
enters the discussion only when a particular theory is named, and it
enters downstream of, not upstream of, (i) and (ii).

The inspection is verified mechanically. In the formal development
(Appendix~\ref{app:lean}, Certificate~4), the master theorem is
quantified over an arbitrary type $G$, and $G$ occurs in none of the
hypotheses, definitions, or conclusions; the proof term is constant
in $G$, witnessed by the fact that the identity of proofs at two
distinct instantiations is definitional. This is the machine-checked
form of claim (iii).
\end{proof}

%%=====================================================================
\section{Scope and outlook}\label{sec:scope}
%%=====================================================================

\begin{corollary}[Scope: the value of $\Delta$ is not gauge-group
independent]\label{cor:scope}
Proposition~\ref{prop:transfer} establishes the gauge-group
independence of the existence claim $\Delta > 0$ only. The numerical
value of $\Delta$ depends on which corner is distinguished and on the
induced metric data --- the selection $(g, \iota, c)$ of
Definition~\ref{def:setting} --- and is therefore not a gauge-group
invariant; no such claim is made here.
\end{corollary}

\begin{proof}
The existence half is uniform: positivity consumes only distinctness
(Lemma~\ref{lem:corners}) and the metric-space structure
(Lemma~\ref{lem:positivity}), and these hold for every admissible
selection. The value half is witnessed by exhibiting two admissible
selections with different gap values; an explicit pair, with gaps in
the ratio two to one, is constructed in the formal development
(Appendix~\ref{app:lean}, Certificate~5). Hence the value is a
non-constant function of the selection data and cannot be an
invariant of anything that does not determine the selection.
\end{proof}

The two halves of Corollary~\ref{cor:scope} have different logical
characters, and the distinction is load-bearing. Existence is
uniform over selections and over gauge groups; value is a function of
the selection. A reader interested in a particular reduced theory
should therefore expect the present results to supply the existence
of its gap, conditional on membership of $\class$, and to be silent
on the gap's magnitude. Which corner is distinguished is likewise
data from the present standpoint; we note only that in reductions
whose internal space carries structure beyond the metric --- an
orientation of the fibre, or a choice of spin or line-bundle data ---
the selection can be induced by that structure rather than left free,
and the value of $\Delta$ is then determined by the same data that
determine the theory.

Two directions merit note. First, the class $\class$ is not empty
as defined: a spectral-data triple with $m(n) = \dodd(n)$ and
$\mathrm{mass}(n) = n\Delta$, over any admissible setting of
Definition~\ref{def:setting}, satisfies (A1)--(A3) by construction,
and the formal development of Appendix~\ref{app:lean} instantiates
exactly such data. The open question is therefore not non-emptiness
but dynamical realisation: whether a reduction of physical interest
produces this spectral data from its geometry. Exhibiting one, with
the pillowcase realised as an orbifold fibre of the internal space
and the degeneracies computed from its geometry, would convert the
conditional results into statements about that theory. The natural
shape of such a candidate is visible in the structure of the count
itself: we observe that anti-periodic boundary conditions on an
order-two quotient produce half-integer moding, that is, spectra
supported on odd multiples of a basic spacing, and
$\dodd(m) = \#\{(n,a) : n(2a+1) = m\}$ is precisely the enumeration
of a level $m$ by such odd multiples
(Proposition~\ref{prop:appell}); the generating-function context of
Section~\ref{sec:moduli} indicates where degeneracies of this form
arise. Second, the structure of
Proposition~\ref{prop:transfer} --- a gap proved from fixed-point
topology and metric completeness, with the gauge group entering
nowhere --- suggests that existence questions for gaps may in general
be separable from the dynamical questions that determine values. An
existence-only result of this kind may be useful in making progress
towards the mass-gap problem for four-dimensional gauge theories,
where it is the existence, not the value, that is asked for.

\subsection*{Acknowledgements}

The author used Claude (Anthropic) and Gemini (Google DeepMind) as
mathematical development and analytical tools during the preparation
of this work. All intellectual content, theoretical framework, and
scientific claims are the sole responsibility of the author.

%%=====================================================================
\appendix
\section{Formal certificates}\label{app:lean}
%%=====================================================================

Seven machine-checked certificates accompany the paper, developed in
Lean~4 (toolchain 4.15.0) against Mathlib (tag v4.15.0). The suite
compiles with zero errors, zero warnings, and no incomplete proofs;
an axiom audit confirms that every headline theorem depends only on
the three standard axioms of the Lean/Mathlib foundation
(propositional extensionality, choice, and quotient soundness).
Certificates 1 and 2 are unconditional and cover
Lemmas~\ref{lem:corners} and~\ref{lem:positivity}, the latter at
metric-space level as noted in Section~\ref{sec:setting}.
Certificate~3 is conditional and machine-checks
Proposition~\ref{prop:identification}: the hypotheses (A1)--(A3), and
exactly those, are encoded as named fields, with the internal input
$\Delta \ge 0$ an explicit argument. Certificate~4 machine-checks
Proposition~\ref{prop:transfer}, with the gauge group a phantom type
parameter and the constancy of the proof term in that parameter
witnessed definitionally. Certificate~5 is the statement-level form
of Corollary~\ref{cor:scope}, including the explicit two-selection
witness. Certificate~6 proves the parity implication of
Theorem~\ref{thm:betti} in full generality, and Certificate~7 proves
the coefficient identity of Proposition~\ref{prop:appell} for all
$m \ge 1$ via the explicit divisor bijection, together with the
certified table to order fifteen.

\subsection*{Certificate 1 --- corner distinctness (\texttt{Corners.lean})}
\begin{Verbatim}[fontsize=\footnotesize,breaklines=true,breakanywhere=true,formatcom=\leanmono]
/-
Certificate 1 (Lemma: corner distinctness).
On the torus T² = ℝ²/ℤ², the involution σ : (x,y) ↦ (-x,-y) has
exactly four fixed points -- the 2-torsion points
(0,0), (1/2,0), (0,1/2), (1/2,1/2) -- and they are pairwise distinct.
Unconditional; zero `sorry`. No reference to any gauge group, and no
quaternionic structure, appears in this file.
Model: the circle is Mathlib's `AddCircle (1 : ℝ) = ℝ/ℤ`; the torus
is the product of two copies -- the object of the prose statement
itself, not a substitute model.
-/
import Mathlib.Topology.Instances.AddCircle

noncomputable section
namespace Corners

open AddCircle

/-- The circle ℝ/ℤ. -/
abbrev S : Type := AddCircle (1 : ℝ)

/-- The torus T² = (ℝ/ℤ)². -/
abbrev T : Type := S × S

/-- The involution σ(x,y) = (-x,-y). -/
def sigma (q : T) : T := (-q.1, -q.2)

/-- The four corner points. -/
def c00 : T := (((0 : ℝ) : S), ((0 : ℝ) : S))
def c10 : T := ((((1 : ℝ)/2 : ℝ) : S), ((0 : ℝ) : S))
def c01 : T := (((0 : ℝ) : S), (((1 : ℝ)/2 : ℝ) : S))
def c11 : T := ((((1 : ℝ)/2 : ℝ) : S), (((1 : ℝ)/2 : ℝ) : S))

/-! ### Single-circle facts -/

/-- An integer is zero on the circle. -/
theorem int_coe_eq_zero (m : ℤ) : (((m : ℝ)) : S) = ((0 : ℝ) : S) := by
  have h : (((m : ℝ)) : S) = 0 :=
    (AddCircle.coe_eq_zero_iff (1 : ℝ)).mpr ⟨m, by rw [zsmul_eq_mul, mul_one]⟩
  rw [h, show ((0 : ℝ) : S) = 0 by simp]

/-- 1/2 is not zero on the circle. -/
theorem half_ne_zero : (((1 : ℝ)/2 : ℝ) : S) ≠ ((0 : ℝ) : S) := by
  intro h
  have h0 : (((1 : ℝ)/2 : ℝ) : S) = 0 := by simpa using h
  rw [AddCircle.coe_eq_zero_iff] at h0
  obtain ⟨n, hn⟩ := h0
  rw [zsmul_eq_mul, mul_one] at hn
  have h2 : ((2 * n : ℤ) : ℝ) = 1 := by push_cast; linarith
  have h3 : (2 * n : ℤ) = 1 := by exact_mod_cast h2
  omega

/-- A point of the circle is fixed by negation iff it is 0 or 1/2. -/
theorem neg_eq_self_iff (x : S) :
    -x = x ↔ x = ((0 : ℝ) : S) ∨ x = (((1 : ℝ)/2 : ℝ) : S) := by
  refine QuotientAddGroup.induction_on x (fun r => ?_)
  constructor
  · intro h
    have hneg : ((-r : ℝ) : S) = ((r : ℝ) : S) := by
      rw [AddCircle.coe_neg]; exact h
    have h0 : ((r + r : ℝ) : S) = 0 := by
      have := sub_eq_zero.mpr hneg.symm
      rw [← AddCircle.coe_sub] at this
      simpa [sub_neg_eq_add] using this
    rw [AddCircle.coe_eq_zero_iff] at h0
    obtain ⟨n, hn⟩ := h0
    rw [zsmul_eq_mul, mul_one] at hn
    rcases Int.even_or_odd n with ⟨m, hm⟩ | ⟨m, hm⟩
    · left
      have hr : r = (m : ℝ) := by
        have : ((m : ℝ) + (m : ℝ)) = r + r := by
          rw [← hn, hm]; push_cast; ring
        linarith
      rw [hr]; exact int_coe_eq_zero m
    · right
      have hr : r = (m : ℝ) + 1/2 := by
        have : ((2 * m + 1 : ℤ) : ℝ) = r + r := by rw [← hn, hm]
        push_cast at this; linarith
      have : ((r - 1/2 : ℝ) : S) = ((0 : ℝ) : S) := by
        rw [show (r - 1/2 : ℝ) = (m : ℝ) by rw [hr]; ring]
        exact int_coe_eq_zero m
      have h5 : ((r : ℝ) : S) - (((1:ℝ)/2 : ℝ) : S) = 0 := by
        rw [← AddCircle.coe_sub]; simpa using this
      exact sub_eq_zero.mp h5
  · rintro (h | h)
    · rw [h]; simp
    · rw [h]
      have : ((-((1:ℝ)/2) : ℝ) : S) = (((1:ℝ)/2 : ℝ) : S) := by
        have : ((-((1:ℝ)/2) - (1:ℝ)/2 : ℝ) : S) = ((0 : ℝ) : S) := by
          rw [show (-((1:ℝ)/2) - (1:ℝ)/2 : ℝ) = ((-1 : ℤ) : ℝ) by push_cast; ring]
          exact int_coe_eq_zero (-1)
        have h5 : ((-((1:ℝ)/2) : ℝ) : S) - (((1:ℝ)/2 : ℝ) : S) = 0 := by
          rw [← AddCircle.coe_sub]; simpa using this
        exact sub_eq_zero.mp h5
      rw [← AddCircle.coe_neg] at *
      exact this

/-! ### The torus theorems -/

/-- **Classification.** The fixed-point set of σ is exactly the four
corner points. -/
theorem fixed_iff (q : T) :
    sigma q = q ↔ q = c00 ∨ q = c10 ∨ q = c01 ∨ q = c11 := by
  obtain ⟨x, y⟩ := q
  have hx := neg_eq_self_iff x
  have hy := neg_eq_self_iff y
  simp only [sigma, Prod.mk.injEq, c00, c10, c01, c11]
  constructor
  · rintro ⟨h1, h2⟩
    rcases hx.mp h1 with rfl | rfl <;> rcases hy.mp h2 with rfl | rfl <;> tauto
  · rintro (⟨rfl, rfl⟩ | ⟨rfl, rfl⟩ | ⟨rfl, rfl⟩ | ⟨rfl, rfl⟩) <;>
      exact ⟨hx.mpr (by tauto), hy.mpr (by tauto)⟩

/-- Each corner is fixed. -/
theorem corners_fixed :
    sigma c00 = c00 ∧ sigma c10 = c10 ∧ sigma c01 = c01 ∧ sigma c11 = c11 :=
  ⟨(fixed_iff c00).mpr (by tauto), (fixed_iff c10).mpr (by tauto),
   (fixed_iff c01).mpr (by tauto), (fixed_iff c11).mpr (by tauto)⟩

/-- **Distinctness.** The four corners are pairwise distinct. -/
theorem corners_pairwise_distinct :
    c00 ≠ c10 ∧ c00 ≠ c01 ∧ c00 ≠ c11 ∧
    c10 ≠ c01 ∧ c10 ≠ c11 ∧ c01 ≠ c11 :=
  ⟨fun h => half_ne_zero (congrArg Prod.fst h).symm,
   fun h => half_ne_zero (congrArg Prod.snd h).symm,
   fun h => half_ne_zero (congrArg Prod.fst h).symm,
   fun h => half_ne_zero (congrArg Prod.fst h),
   fun h => half_ne_zero (congrArg Prod.snd h).symm,
   fun h => half_ne_zero (congrArg Prod.fst h).symm⟩

/-- **Summary.** σ has exactly the four corners as fixed points, and
they are pairwise distinct. -/
theorem corner_distinctness :
    (∀ q : T, sigma q = q ↔ q = c00 ∨ q = c10 ∨ q = c01 ∨ q = c11) ∧
    (c00 ≠ c10 ∧ c00 ≠ c01 ∧ c00 ≠ c11 ∧
     c10 ≠ c01 ∧ c10 ≠ c11 ∧ c01 ≠ c11) :=
  ⟨fixed_iff, corners_pairwise_distinct⟩

end Corners
\end{Verbatim}

\subsection*{Certificate 2 --- positivity of the corner gap (\texttt{Gap.lean})}
\begin{Verbatim}[fontsize=\footnotesize,breaklines=true,breakanywhere=true,formatcom=\leanmono]
/-
Certificate 2 (Lemma: positivity of the corner gap).
Distinct points of a metric space are at strictly positive distance;
instantiated at the four corner points of Certificate 1, giving Δ > 0.
Unconditional; zero `sorry`.
Formalisation note (stated once): the certificate is proved at
METRIC-SPACE level -- strictly weaker machinery than Hopf--Rinow and
strictly sufficient for the claim as used. The prose retains
Hopf--Rinow as the Riemannian-geometric source: the ambient manifold
M, compact with continuous Riemannian metric, induces a metric-space
structure, and `ι` below is the injective inclusion of the pillowcase
fibre carrying the induced distance.
No reference to any gauge group appears in this file.
-/
import YML.Corners

noncomputable section
namespace Gap
open Corners

/-- Distinct points of a metric space are at strictly positive
distance. -/
theorem dist_pos_of_ne {X : Type*} [MetricSpace X] {a b : X}
    (h : a ≠ b) : 0 < dist a b :=
  dist_pos.mpr h

/-- The minimum of three positive reals is positive. -/
theorem min₃_pos {a b c : ℝ} (ha : 0 < a) (hb : 0 < b) (hc : 0 < c) :
    0 < min a (min b c) :=
  lt_min ha (lt_min hb hc)

section Embedding

variable {K : Type*} [MetricSpace K]
variable (ι : T → K)

/-- All six pairwise corner distances in `K` are positive. -/
theorem corner_dists_pos (hι : Function.Injective ι) :
    0 < dist (ι c00) (ι c10) ∧ 0 < dist (ι c00) (ι c01) ∧
    0 < dist (ι c00) (ι c11) ∧ 0 < dist (ι c10) (ι c01) ∧
    0 < dist (ι c10) (ι c11) ∧ 0 < dist (ι c01) (ι c11) := by
  obtain ⟨h1, h2, h3, h4, h5, h6⟩ := corners_pairwise_distinct
  exact ⟨dist_pos_of_ne (fun h => h1 (hι h)),
         dist_pos_of_ne (fun h => h2 (hι h)),
         dist_pos_of_ne (fun h => h3 (hι h)),
         dist_pos_of_ne (fun h => h4 (hι h)),
         dist_pos_of_ne (fun h => h5 (hι h)),
         dist_pos_of_ne (fun h => h6 (hι h))⟩

/-- **Positivity of the corner gap.** The minimum distance from the
distinguished corner to the three remaining corners is strictly
positive. Stated with `c00` as the distinguished exemplar; the
argument is corner-symmetric, and the selection data pick which
corner plays the role. -/
theorem gap_pos (hι : Function.Injective ι) :
    0 < min (dist (ι c00) (ι c10))
        (min (dist (ι c00) (ι c01)) (dist (ι c00) (ι c11))) := by
  obtain ⟨h1, h2, h3, _, _, _⟩ := corner_dists_pos ι hι
  exact min₃_pos h1 h2 h3

end Embedding
end Gap
\end{Verbatim}

\subsection*{Certificate 3 --- conditional identification (\texttt{Identification.lean})}
\begin{Verbatim}[fontsize=\footnotesize,breaklines=true,breakanywhere=true,formatcom=\leanmono]
/-
Certificate 3 (Proposition: conditional identification).
Δ is the mass of the lowest non-vacuum one-particle state of a theory
in the class C. CONDITIONAL: the certificate machine-checks the
conditional structure -- the three class hypotheses (A1)-(A3), and
exactly those three, are encoded as named fields, and the conclusion
is derived from them. It does not prove the hypotheses.
  hA1 -- (A1): level degeneracies m(n) = d_odd(n)
  hA2 -- (A2): the mass at level n equals n·Δ
  hA3 -- (A3): d_odd(n) counts one-particle states at level n, so the
        non-vacuum one-particle masses are exactly the occupied
        levels n ≥ 1
d_odd itself is NOT a hypothesis: it is defined arithmetically (count
of odd divisors) and its needed value d_odd(1) = 1 is proved.
FINDING (made visible rather than silent): the lower-bound direction
needs 0 ≤ Δ. In the prose this is supplied by the positivity lemma;
here it is an explicit argument `hΔpos`, and the master certificate
(Certificate 4) discharges it from the geometry.
No reference to any gauge group appears in this file.
-/
import Mathlib.NumberTheory.Divisors
import Mathlib.Analysis.SpecialFunctions.Pow.Real

noncomputable section
namespace Identification

/-- The odd-divisor count d_odd(n), defined arithmetically. -/
def dodd (n : ℕ) : ℕ := ((Nat.divisors n).filter (fun d => Odd d)).card

/-- The smallest case, worked by machine: d_odd(1) = 1. -/
theorem dodd_one : dodd 1 = 1 := by decide

/-- The identification hypotheses, exactly the three of the class. -/
structure IdentificationData where
  /-- the internal spectral separation between the distinguished
      corner and the nearest remaining corner (geometric input;
      positive by the gap lemma when instantiated -- assumed here,
      discharged in Certificate 4) -/
  Δ : ℝ
  /-- level degeneracies of the reduced theory -/
  m : ℕ → ℕ
  /-- (A1): m(n) = d_odd(n) -/
  hA1 : ∀ n, m n = dodd n
  /-- one-particle mass assignment at level n -/
  mass : ℕ → ℝ
  /-- (A2): the activated spectrum is Δ-spaced -/
  hA2 : ∀ n, mass n = n * Δ
  /-- the set of masses of non-vacuum one-particle states -/
  massSet : Set ℝ
  /-- (A3): the non-vacuum one-particle masses are exactly the
      occupied levels n ≥ 1 -/
  hA3 : massSet = { x | ∃ n : ℕ, 1 ≤ n ∧ m n ≠ 0 ∧ x = mass n }

/-- **Conditional identification.** Under the three hypotheses, Δ is
the least element of the set of non-vacuum one-particle masses. -/
theorem identification (D : IdentificationData) (hΔpos : 0 ≤ D.Δ) :
    IsLeast D.massSet D.Δ := by
  constructor
  · -- membership: level 1 is occupied (d_odd(1) = 1) and has mass Δ
    rw [D.hA3]
    refine ⟨1, le_refl 1, ?_, ?_⟩
    · rw [D.hA1, dodd_one]; exact one_ne_zero
    · rw [D.hA2]; simp
  · -- lower bound: every occupied level n ≥ 1 has mass nΔ ≥ Δ
    rintro x hx
    rw [D.hA3] at hx
    obtain ⟨n, h1, _, rfl⟩ := hx
    rw [D.hA2]
    calc D.Δ = 1 * D.Δ := (one_mul _).symm
    _ ≤ n * D.Δ := by
        apply mul_le_mul_of_nonneg_right _ hΔpos
        exact_mod_cast h1

end Identification
\end{Verbatim}

\subsection*{Certificate 4 --- gauge transfer (\texttt{Transfer.lean})}
\begin{Verbatim}[fontsize=\footnotesize,breaklines=true,breakanywhere=true,formatcom=\leanmono]
/-
Certificate 4 (Proposition: gauge transfer, master form).
The transfer proposition assembled, with claim (iii) -- gauge-group
independence -- made STRUCTURAL. The master theorem is quantified over
an arbitrary type `G` ("the gauge group"), and `G` occurs in NONE of
the hypotheses, definitions, or conclusions: the proof term is
constant in `G` (verified by `rfl`). This is the machine-checkable
form of the inspection of the proof's dependencies: the five inputs --
  (1) the 2-torsion structure of σ on T²      [Certificate 1, proved]
  (2) positive distance between distinct points [Certificate 2, proved]
  (3) (A1), (4) (A2), (5) (A3)                 [Certificate 3, named]
-- are all G-free, verified by Lean's binder discipline.
Conditional on (A1)-(A3); unconditional in G.
-/
import YML.Corners
import YML.Gap
import YML.Identification

noncomputable section
namespace Transfer
open Corners Gap Identification

section Master

variable {K : Type*} [MetricSpace K]   -- the ambient manifold, as a metric space
variable (ι : T → K)                    -- inclusion of the pillowcase fibre

/-- The corner gap: distance from the distinguished corner to the
nearest remaining corner. `c00` as the distinguished exemplar; the
argument is corner-symmetric. -/
def cornerGap : ℝ :=
  min (dist (ι c00) (ι c10)) (min (dist (ι c00) (ι c01)) (dist (ι c00) (ι c11)))

-- NOTE: without the option below, Lean's linter reports
-- "unused variable `G`" on the theorem that follows. That report IS
-- claim (iii), issued by the machine itself. It is silenced only to
-- keep the compile warning-free, having been placed on the record.
set_option linter.unusedVariables false in
/-- **Gauge transfer (master form).** For an ARBITRARY gauge group
`G` -- which occurs nowhere in the hypotheses -- if the fibre embeds
injectively in `K` and the identification data carries Δ = cornerGap,
then (i) Δ > 0, and (ii) Δ is the mass of the lowest non-vacuum
one-particle state. Claim (iii) is the visible fact that `G` is a
phantom parameter. -/
theorem gauge_transfer (G : Type*)
    (hι : Function.Injective ι)
    (D : IdentificationData)
    (hΔdef : D.Δ = cornerGap ι) :
    0 < D.Δ ∧ IsLeast D.massSet D.Δ := by
  have hpos : 0 < D.Δ := by rw [hΔdef]; exact gap_pos ι hι
  exact ⟨hpos, identification D (le_of_lt hpos)⟩

/-- The proof term is literally constant in `G`. -/
theorem transfer_constant_in_G (G₁ G₂ : Type*)
    (hι : Function.Injective ι) (D : IdentificationData)
    (hΔdef : D.Δ = cornerGap ι) :
    gauge_transfer ι G₁ hι D hΔdef = gauge_transfer ι G₂ hι D hΔdef :=
  rfl

end Master
end Transfer
\end{Verbatim}

\subsection*{Certificate 5 --- scope (\texttt{Scope.lean})}
\begin{Verbatim}[fontsize=\footnotesize,breaklines=true,breakanywhere=true,formatcom=\leanmono]
/-
Certificate 5 (Corollary: scope, statement-level).
The EXISTENCE claim Δ > 0 is gauge-group independent; the numerical
VALUE of Δ is not: it depends on the corner-selection data, and no
value-independence claim is made.
Formalisation (statement-level): a "selection" is the data the value
actually depends on -- the distinguished corner image and the three
remaining corner images, pairwise distinct from it. Two theorems:
  (1) existence is selection- and G-uniform;
  (2) two explicit admissible selections with DIFFERENT gap values
      (1 and 2), witnessing that the value is a function of the
      selection, not a gauge-group invariant.
As in Certificate 4, `G` is a phantom parameter throughout.
Unconditional (statement-level); zero `sorry`.
-/
import YML.Gap

noncomputable section
namespace Scope
open Gap

/-- A corner selection in a metric space `K`: the distinguished corner
image and the three remaining corner images, pairwise distinct from
it. (The data the value of Δ actually depends on.) -/
structure Selection (K : Type*) [MetricSpace K] where
  vac : K
  o1 : K
  o2 : K
  o3 : K
  h1 : vac ≠ o1
  h2 : vac ≠ o2
  h3 : vac ≠ o3

/-- The gap of a selection. -/
def Selection.gap {K : Type*} [MetricSpace K] (s : Selection K) : ℝ :=
  min (dist s.vac s.o1) (min (dist s.vac s.o2) (dist s.vac s.o3))

-- The linter's "unused variable `G`" report is the scope claim's
-- gauge-independence half, issued by the machine; silenced only to
-- keep the compile warning-free, having been placed on the record.
set_option linter.unusedVariables false in
/-- **Existence half.** For every admissible selection and every
gauge group `G`, the gap is positive. Selection- and G-uniform. -/
theorem existence_selection_and_gauge_independent
    (G : Type*) {K : Type*} [MetricSpace K] (s : Selection K) :
    0 < s.gap :=
  min₃_pos (dist_pos_of_ne s.h1) (dist_pos_of_ne s.h2) (dist_pos_of_ne s.h3)

/-! ### Value half: two admissible selections, different values -/

/-- Selection A: corners at 0, 1, 2, 3 on the real line. -/
def selA : Selection ℝ where
  vac := 0
  o1 := 1
  o2 := 2
  o3 := 3
  h1 := by norm_num
  h2 := by norm_num
  h3 := by norm_num

/-- Selection B: corners at 0, 2, 4, 6 on the real line. -/
def selB : Selection ℝ where
  vac := 0
  o1 := 2
  o2 := 4
  o3 := 6
  h1 := by norm_num
  h2 := by norm_num
  h3 := by norm_num

theorem gapA : selA.gap = 1 := by
  simp only [Selection.gap, selA, Real.dist_eq]
  norm_num

theorem gapB : selB.gap = 2 := by
  simp only [Selection.gap, selB, Real.dist_eq]
  norm_num

/-- **Value half.** The value of the gap is not
selection-independent: two admissible selections with different
values. Hence no value-independence claim can be made. -/
theorem value_depends_on_selection : selA.gap ≠ selB.gap := by
  rw [gapA, gapB]; norm_num

end Scope
\end{Verbatim}

\subsection*{Certificate 6 --- signed count equals absolute count (\texttt{BettiParity.lean})}
\begin{Verbatim}[fontsize=\footnotesize,breaklines=true,breakanywhere=true,formatcom=\leanmono]
/-
Certificate 6 (Theorem: signed count equals absolute count).
The arithmetic implication of the index theorem: if all odd Betti
numbers vanish, the signed Euler characteristic equals the absolute
Betti count -- there are no cancellations. The geometric inputs (the
cellular decomposition and the vanishing itself) are external
theorems cited in the prose; this certificate machine-checks the
implication, in full generality (any Betti sequence, any range),
not merely at statement level.
Unconditional; zero `sorry`.
-/
import Mathlib.Algebra.BigOperators.Group.Finset
import Mathlib.Tactic.NormNum
import Mathlib.Algebra.Order.BigOperators.Group.Finset

namespace BettiParity

open Finset

/-- **Signed = absolute.** For any Betti sequence b : ℕ → ℤ with
vanishing odd part, the alternating sum equals the plain sum over any
initial range. -/
theorem signed_eq_absolute (b : ℕ → ℤ)
    (hodd : ∀ j, b (2 * j + 1) = 0) (N : ℕ) :
    ∑ j ∈ range N, (-1 : ℤ) ^ j * b j = ∑ j ∈ range N, b j := by
  apply Finset.sum_congr rfl
  intro j _
  rcases Nat.even_or_odd j with ⟨k, hk⟩ | ⟨k, hk⟩
  · -- even j: the sign is +1
    subst hk
    rw [show k + k = 2 * k by omega, pow_mul]
    norm_num
  · -- odd j: the term vanishes on both sides
    subst hk
    rw [hodd k, mul_zero]

/-- Corollary in the form used in the prose: under odd vanishing, the
Euler characteristic (alternating sum) is a sum of non-negative terms
whenever the Betti numbers are non-negative. -/
theorem chi_nonneg (b : ℕ → ℤ) (hodd : ∀ j, b (2 * j + 1) = 0)
    (hpos : ∀ j, 0 ≤ b j) (N : ℕ) :
    0 ≤ ∑ j ∈ range N, (-1 : ℤ) ^ j * b j := by
  rw [signed_eq_absolute b hodd N]
  exact Finset.sum_nonneg (fun j _ => hpos j)

end BettiParity
\end{Verbatim}

\subsection*{Certificate 7 --- Appell--Lerch central coefficients (\texttt{OddDivisors.lean})}
\begin{Verbatim}[fontsize=\footnotesize,breaklines=true,breakanywhere=true,formatcom=\leanmono]
/-
Certificate 7 (Proposition: Appell--Lerch central coefficients).
The central (y⁰) coefficient of the level-two Appell--Lerch sum
μ(τ,z) = Σ_{n≥1} qⁿ/((1-yqⁿ)(1-y⁻¹qⁿ)) at qᵐ counts the pairs
(n,a) with n ≥ 1, a ≥ 0 and n(2a+1) = m -- equivalently, the odd
divisors of m. This certificate machine-checks (i) the identity
pairCount m = d_odd m for ALL m ≥ 1, via the explicit bijection
d = 2a+1 with inverse (n,a) = (m/d, (d-1)/2), and (ii) the
coefficient table to q¹⁵ by decidable evaluation. The
geometric-series expansion step from μ to pairCount is a two-line
manipulation carried in the prose.
Unconditional; zero `sorry`.
-/
import Mathlib.NumberTheory.Divisors

namespace OddDivisors

open Finset

/-- The odd-divisor count d_odd(m). -/
def dodd (m : ℕ) : ℕ := ((Nat.divisors m).filter (fun d => Odd d)).card

/-- The coefficient count: pairs (n,a), n ≥ 1, a ≥ 0, with
n(2a+1) = m. Ranges suffice since n ≤ m and (d-1)/2 ≤ m for m ≥ 1. -/
def pairCount (m : ℕ) : ℕ :=
  (((range (m + 1)) ×ˢ (range (m + 1))).filter
    (fun p => 1 ≤ p.1 ∧ p.1 * (2 * p.2 + 1) = m)).card

/-- **The identity, in general.** For all m ≥ 1, the coefficient count
equals the odd-divisor count. -/
theorem pairCount_eq_dodd (m : ℕ) (hm : 1 ≤ m) :
    pairCount m = dodd m := by
  unfold pairCount dodd
  apply Finset.card_nbij' (fun p => 2 * p.2 + 1)
    (fun d => (m / d, (d - 1) / 2))
  · -- forward map lands in the odd divisors
    rintro ⟨n, a⟩ hp
    simp only [mem_filter, mem_product, mem_range] at hp
    obtain ⟨⟨-, -⟩, hn1, hprod⟩ := hp
    simp only [mem_filter, Nat.mem_divisors]
    exact ⟨⟨⟨n, by rw [Nat.mul_comm]; exact hprod.symm⟩, by omega⟩,
           ⟨a, rfl⟩⟩
  · -- inverse map lands in the pair set
    intro d hd
    simp only [mem_filter, Nat.mem_divisors] at hd
    obtain ⟨⟨hdvd, hm0⟩, hodd⟩ := hd
    have hd1 : 1 ≤ d := Nat.one_le_iff_ne_zero.mpr (by rintro rfl; simp at hodd)
    have hdle : d ≤ m := Nat.le_of_dvd (by omega) hdvd
    have hq1 : 1 ≤ m / d := Nat.one_le_div_iff (by omega) |>.mpr hdle
    have hqle : m / d ≤ m := Nat.div_le_self m d
    have hdodd : d % 2 = 1 := Nat.odd_iff.mp hodd
    simp only [mem_filter, mem_product, mem_range]
    refine ⟨⟨by omega, by omega⟩, hq1, ?_⟩
    rw [show 2 * ((d - 1) / 2) + 1 = d by omega]
    exact Nat.div_mul_cancel hdvd
  · -- left inverse
    rintro ⟨n, a⟩ hp
    simp only [mem_filter, mem_product, mem_range] at hp
    obtain ⟨⟨-, -⟩, hn1, hprod⟩ := hp
    have hpos : 0 < 2 * a + 1 := by omega
    have hdiv : m / (2 * a + 1) = n := by
      rw [← hprod, Nat.mul_div_cancel _ hpos]
    simp only [Prod.mk.injEq]
    exact ⟨hdiv, by omega⟩
  · -- right inverse
    intro d hd
    simp only [mem_filter, Nat.mem_divisors] at hd
    have hdodd : d % 2 = 1 := Nat.odd_iff.mp hd.2
    have hd1 : 1 ≤ d := by omega
    dsimp only
    omega

/-- **The table to q¹⁵**, by decidable evaluation. -/
theorem table :
    (List.range 15).map (fun i => dodd (i + 1)) =
    [1, 1, 2, 1, 2, 2, 2, 1, 3, 2, 2, 2, 2, 2, 4] := by decide

/-- The smallest case, restated: d_odd(1) = 1. -/
theorem dodd_one : dodd 1 = 1 := by decide

/-- Strict positivity to q¹⁵, by decidable evaluation, accompanying
the one-line general proof in the prose. -/
theorem positive_to_fifteen : ∀ m ∈ List.range 15, 1 ≤ dodd (m + 1) := by
  decide

end OddDivisors
\end{Verbatim}


\begin{thebibliography}{9}

\bibitem{BB}
A.~Bia{\l}ynicki-Birula,
Some theorems on actions of algebraic groups,
\textit{Ann.\ of Math.}\ \textbf{98} (1973), 480--497.

\bibitem{BMR}
K.~Bringmann, J.~Manschot and L.~Rolen,
Identities for generalized Appell functions and the blow-up formula,
\textit{Lett.\ Math.\ Phys.}\ \textbf{106} (2016), 1379--1395.

\bibitem{doCarmo}
M.~P.~do Carmo,
\textit{Riemannian Geometry},
Birkh\"auser, Boston, 1992.

\bibitem{Gottsche}
L.~G\"ottsche,
The Betti numbers of the Hilbert scheme of points on a smooth
projective surface,
\textit{Math.\ Ann.}\ \textbf{286} (1990), 193--207.

\bibitem{NakajimaLectures}
H.~Nakajima,
\textit{Lectures on Hilbert Schemes of Points on Surfaces},
University Lecture Series \textbf{18}, American Mathematical Society,
Providence, RI, 1999.

\bibitem{NakajimaYoshioka}
H.~Nakajima and K.~Yoshioka,
Instanton counting on blowup.\ I.\ 4-dimensional pure gauge theory,
\textit{Invent.\ Math.}\ \textbf{162} (2005), 313--355.

\bibitem{NYlectures}
H.~Nakajima and K.~Yoshioka,
Lectures on instanton counting,
in \textit{Algebraic Structures and Moduli Spaces},
CRM Proc.\ Lecture Notes \textbf{38}, American Mathematical Society,
Providence, RI, 2004, 31--101.

\bibitem{Serre}
J.-P.~Serre,
\textit{A Course in Arithmetic},
Graduate Texts in Mathematics \textbf{7}, Springer, New York, 1973.

\bibitem{VafaWitten}
C.~Vafa and E.~Witten,
A strong coupling test of $S$-duality,
\textit{Nucl.\ Phys.\ B} \textbf{431} (1994), 3--77.

\bibitem{Zwegers}
S.~Zwegers,
\textit{Mock Theta Functions},
Ph.D.\ thesis, Universiteit Utrecht, 2002.

\end{thebibliography}

\end{document}
